% Biểu tượng tiêu đề được thay đổi theo từng phần.
\newcommand{\subsectionicon}{}

% Bắt đầu mỗi bài học trên trang mới và đánh số trang lại từ 1.
% (chia sẽ từng bài trước khi thành sách)
\newcommand{\sectionbreak}{\clearpage\thispagestyle{sectionfirst}\pagenumbering{arabic}}

% Tiêu đề bài học: số ở bên trái, tên bài ở bên phải.
\titlespacing{\section}{0pt}{0pt}{0pt}
\titleformat{\section}[block]
{}
{}
{0pt}
{%
    \begin{tikzpicture}[remember picture, overlay]
        \begin{scope}[shift={([xshift=48pt]current page.north west)}]
            \path
            (0,0) coordinate (NW)
            (56pt,0) coordinate (NE)
            (0,-48pt) coordinate (SW)
            (56pt,-48pt) coordinate (SE)
            (28pt,-72pt) coordinate (TIP);
            \fill[base700] (NW) -- (NE) -- (SE) -- (TIP) -- (SW) -- cycle;
            \path[name path=left-angled-helper] ($(SW)!4pt!90:(TIP)$) -- ($(TIP)!4pt!-90:(SW)$);
            \path[name path=right-angled-helper] ($(SE)!4pt!-90:(TIP)$) -- ($(TIP)!4pt!90:(SE)$);
            \path[name intersections={of=left-angled-helper and right-angled-helper, by=stitch-tip}];
            \path[name path=left-vertical-helper] ([xshift=4pt]NW) -- ([xshift=4pt]SW);
            \path[name path=right-vertical-helper] ([xshift=-4pt]NE) -- ([xshift=-4pt]SE);
            \path[name intersections={of=left-vertical-helper and left-angled-helper, by=inner-sw}];
            \path[name intersections={of=right-vertical-helper and right-angled-helper, by=inner-se}];
            \draw[white, dashed, dash pattern=on 3pt off 2pt]
            ([xshift=4pt]NW) -- (inner-sw) -- (stitch-tip) -- (inner-se) -- ([xshift=-4pt]NE);
            \node[text=white, font={\LARGE\titlefont}] at (28pt, -32pt) {\thesection};
        \end{scope}
        \node [
            anchor=north west,
            align=left,
            font=\color{base800}\titlefont\Large,
            minimum height=64pt
        ]
        at ([xshift=136pt, yshift=-8pt]current page.north west)
        {\uppercase{#1}};
    \end{tikzpicture}%
}

% Tiêu đề section không phải bài: không hiển thị số bài và căn giữa tiêu đề.
\titleformat{name=\section,numberless}[block]
{}
{}
{0pt}
{%
    \begin{tikzpicture}[remember picture, overlay]
        \node [
            anchor=north,
            align=center,
            font=\color{base800}\titlefont\Large,
            minimum height=64pt
        ]
        at ([yshift=-8pt]current page.north)
        {\uppercase{#1}};
    \end{tikzpicture}%
}

% Tiêu đề mục có biểu tượng và nền tự co theo nội dung.
\titlespacing{\subsection}{0pt}{8pt}{4pt}
\titleformat{\subsection}[block]
{\vspace{8pt}}
{}
{0pt}
{%
    \begin{tikzpicture}[remember picture, overlay]
        \node [anchor=west, inner ysep=6pt, inner xsep=8pt] (title) at (14pt, 0) {\color{base800}\small\titlefont\uppercase{#1}};
        \path let \p{north} = (title.north), \p{south} = (title.south) in
        \pgfextra{
            \pgfmathsetlengthmacro{\squaresize}{\y{north}-\y{south}}
            \xdef\globalSquaresize{\squaresize}
        }
        node [left color=base700, right color=base600, minimum size=\squaresize, rounded corners=8pt] (icon) {\subsectionicon};
        \begin{scope}[on background layer]
            \path [draw=none, left color=base200, right color=base100]
            ($(icon.south east) + (-8pt, 0)$) --
            (title.south east) --
            ($(title.south east) + (16pt,0.5*\globalSquaresize)$) --
            (title.north east) --
            ($(icon.north east) + (-8pt, 0)$) --
            cycle;
        \end{scope}
    \end{tikzpicture}%
}
[\strut]

% Tiêu đề cấp 2 dùng số La Mã và chữ in hoa.
\titlespacing{\subsubsection}{0pt}{8pt}{4pt}
\titleformat{\subsubsection}[block]
{\small\bfseries}
{\thesubsubsection.}
{0.5em}
{\uppercase{#1}}

% Tiêu đề đoạn dùng số Ả Rập.
\titlespacing{\paragraph}{0pt}{8pt}{4pt}
\titleformat{\paragraph}[block]
{\bfseries}
{\theparagraph.}
{0.5em}
{#1}
